\subsection{Project Deviations}
\label{project-deviations}

With the actual project cost calculated, the deviation formulas defined earlier in Section \ref{sec:management-control} are applied to compute the project deviations across different areas. Table \ref{table:cost-deviation-analysis} below summarizes the cost deviations for human resources, amortization, and total cost. 

\vspace{-10pt}

\normalsize

%Table of hourly wage of each personnel member
\begin{table}[H]
    \renewcommand{\arraystretch}{1.2}
    \centering
    \small
    \setlength{\tabcolsep}{5pt}
    \makebox[\textwidth][c]{%
    \resizebox{0.6\paperwidth}{!}{%
    \begin{tabular}{@{}l r r r@{}}
    \textbf{Cost Category} & \textbf{Planned Cost (\texteuro)} & \textbf{Actual Cost (\texteuro)} & \textbf{Cost Deviation (\texteuro)} \\
    \hline
    Human Resources & 11,913.05 & 13,069.94 & -1,156.89 \\
    Amortization & 102.22 & 108.09 & -5.87 \\
    Total Cost & 14,592.78 & 14,309.00 & 283.78 \\
    \end{tabular}%
    }}
    \normalsize
    \caption[Project cost deviation analysis]
    {Project cost deviation analysis. [Own creation]}
    \vspace{-1em}
    \label{table:cost-deviation-analysis}
\end{table}

\vspace{-10pt}

Although personnel and amortization costs exceeded the initial estimate, contingency and incidental costs covered the budget shortfall. As a result, the overall project cost remained within the allocated budget, thereby demonstrating that the original cost estimation was well-founded.
